% Part: history 
% Chapter: biographies 
% Section: julia-robinson
\documentclass[../../../include/open-logic-section]{subfiles}

\begin{document}

\olfileid{his}{bio}{rob}

\olsection{జూలియా రాబిన్సన్}

\olphoto{robinson-julia}{జూలియా రాబిన్సన్}

జూలియా బౌమన్ రాబిన్సన్ అమెరికన్ గణితశాస్త్రజ్ఞురాలు. నిర్ణయ సమస్యలపై, ముఖ్యంగా హిల్బర్ట్ పదవ సమస్య పరిష్కారానికి ఆమె చేసిన కృషికి ప్రసిద్ధి. రాబిన్సన్ 1919 డిసెంబరు 8న మిస్సౌరీలోని సెయింట్ లూయిస్‌లో జన్మించారు. చిన్నతనం నుంచే సంఖ్యలపై ఆసక్తి కలిగిందని ఆమె గుర్తుచేసుకున్నారు \citep[4]{Reid1986}. తొమ్మిదేళ్ల వయసులో ఆమెకు స్కార్లెట్ జ్వరం సోకింది; రుమాటిక్ జ్వరం మళ్లీ మళ్లీ వచ్చింది. అందువల్ల ఎక్కువ సమయం మంచంపైనే ఉండాల్సి వచ్చి, చదువులో వెనుకబడ్డారు. వ్యక్తిగత బోధకుల సహాయంతో చదువులో తోటి విద్యార్థులను అందుకోగలిగినా, అనారోగ్యం శారీరకంగా ఆమె జీవితంపై దీర్ఘకాలిక ప్రభావం చూపింది.

చిన్నతనంలోని కష్టాలున్నప్పటికీ రాబిన్సన్ గణితం, విజ్ఞానశాస్త్రాలలో అనేక బహుమతులతో ఉన్నత పాఠశాల పూర్తి చేశారు. శాన్ డియెగో స్టేట్ కాలేజీలో విశ్వవిద్యాలయ చదువు మొదలుపెట్టి, చివరి సంవత్సరంలో బర్కిలీలోని కాలిఫోర్నియా విశ్వవిద్యాలయానికి మారారు. అక్కడ గణితశాస్త్రజ్ఞుడు రాఫయెల్ రాబిన్సన్ ప్రభావం ఆమెపై పడింది. వారు మంచి స్నేహితులై 1941లో వివాహం చేసుకున్నారు. అధ్యాపక సభ్యుని జీవిత భాగస్వామి కావడంతో బర్కిలీ గణిత విభాగంలో బోధించడానికి ఆమెకు అనుమతి లేదు. గణిత తరగతులకు శ్రోతగా హాజరవుతూనే ఉన్నా, విశ్వవిద్యాలయాన్ని వదిలి కుటుంబాన్ని ప్రారంభించాలని అనుకున్నారు. కానీ వివాహం జరిగిన కొద్దికాలానికే ఆమెకు న్యుమోనియా సోకింది. చిన్నతనంలో వచ్చిన రుమాటిక్ జ్వరం వల్ల గుండెపై గణనీయమైన మచ్చ కణజాలం ఏర్పడిందని ఆమెకు చెప్పారు. అది తీవ్రంగా ఉండడంతో ఆమె నలభై ఏళ్లు దాటి జీవించరని వైద్యుడు అంచనా వేశారు; పిల్లలను కనవద్దని సూచించారు \citep[13]{Reid1986}.

రాబిన్సన్ చాలాకాలం నిస్పృహతో గడిపారు; కానీ చివరికి గణిత అధ్యయనాన్ని కొనసాగించాలని నిర్ణయించుకున్నారు. బర్కిలీకి తిరిగివెళ్లి ఆల్ఫ్రెడ్ టార్స్కీ పర్యవేక్షణలో 1948లో పీహెచ్‌డీ పూర్తి చేశారు. వాస్తవ సంఖ్యల ప్రథమ శ్రేణి సిద్ధాంతం నిర్ణయించదగినదని టార్స్కీ చూపారు; సహజ సంఖ్యల ప్రథమ శ్రేణి సిద్ధాంతం అనిర్ణయనీయమని గ్యోడెల్ కృషి నుంచి వచ్చింది. పరిమేయ సంఖ్యల ప్రథమ శ్రేణి సిద్ధాంతం నిర్ణయించదగినదా కాదా అనేది ప్రధాన బహిరంగ సమస్యగా ఉండేది. తన పరిశోధన గ్రంథంలో \citeyearpar{Robinson1949} అది నిర్ణయించదగినది కాదని రాబిన్సన్ నిరూపించారు.

నిర్ణయ సమస్యలపై ఆసక్తితో రాబిన్సన్ తరువాత హిల్బర్ట్ పదవ సమస్యను పరిష్కరించడానికి ప్రయత్నించారు. 1900లో డేవిడ్ హిల్బర్ట్ ప్రతిపాదించిన ప్రసిద్ధ 23 గణిత సమస్యల జాబితాలో అది ఒకటి. $3x^2 - 2y +3 = 0$ వంటి పూర్ణసంఖ్య గుణకాలున్న బహుపది సమీకరణానికి పూర్ణసంఖ్యల్లో సాధన ఉందా లేదా అని పరిమిత సమయంలో చెప్పే అల్గోరిథం ఉందా అనేది పదవ సమస్య. అటువంటి ప్రశ్నలను \emph{డయోఫాంటైన్ సమస్యలు} అంటారు. కొన్ని తొలి విజయాల తరువాత, అదే సమస్యపై పనిచేస్తున్న మార్టిన్ డేవిస్, హిలరీ పుట్నమ్‌లతో రాబిన్సన్ కలిశారు. ఘాతీయ డయోఫాంటైన్ సమస్యలు (తెలియని సంఖ్యలు ఘాతాలుగా కూడా కనిపించేవి) అనిర్ణయనీయమని వారు చూపారు; తరువాత ``జె.ఆర్.'' అని పిలిచిన ఒక నిర్దిష్ట పరికల్పన నిజమైతే హిల్బర్ట్ పదవ సమస్య కూడా అనిర్ణయనీయమని చూపారు \citep{DavisPutnamRobinson1961}. 1960లంతా రాబిన్సన్ ఆ సమస్యపై పనిచేశారు. 1970లో యువ రష్యన్ గణితశాస్త్రజ్ఞుడు యూరీ మతియాసేవిచ్ చివరకు జె.ఆర్. పరికల్పనను నిరూపించారు. సంయుక్త ఫలితాన్ని ఇప్పుడు మతియాసేవిచ్--రాబిన్సన్--డేవిస్--పుట్నమ్ సిద్ధాంతం, సంక్షిప్తంగా ఎంఆర్‌డీపీ సిద్ధాంతం అంటారు. మతియాసేవిచ్, రాబిన్సన్ స్నేహితులై అనేక వ్యాసాల్లో కలిసి పనిచేశారు. మతియాసేవిచ్‌కు రాసిన ఒక లేఖలో రాబిన్సన్, ``వేల మైళ్లు దూరంగా ఉన్నా కలిసి పనిచేయడం వల్ల, ఒక్కొక్కరుగా పనిచేసిన దానికంటే స్పష్టంగా ఎక్కువ పురోగతి సాధిస్తున్నందుకు నాకు నిజంగా చాలా సంతోషంగా ఉంది'' అని రాశారు \citep[45]{Matijasevich1992}.

అమెరికన్ మాథమెటికల్ సొసైటీకి అధ్యక్షురాలైన తొలి మహిళ రాబిన్సన్; నేషనల్ అకాడమీ ఆఫ్ సైన్స్‌కు ఎన్నికైన తొలి మహిళ కూడా ఆమె. లుకేమియా నిర్ధారణ అయిన తరువాత, 1985 జూలై 30న 65 ఏళ్ల వయసులో మరణించారు.

\begin{reading}
రాబిన్సన్ గణిత వ్యాసాలు ఆమె \textit{Collected
  Works}లో లభిస్తాయి \citep{Robinson1996}; అందులో ఆమె నేషనల్ అకాడమీ ఆఫ్ సైన్సెస్ జీవిత చరిత్ర స్మృతిరచన \citep{Feferman1994} పునర్ముద్రణ కూడా ఉంది. రాబిన్సన్ అక్క కాన్‌స్టెన్స్ రీడ్, ముఖాముఖిల ఆధారంగా ``Autobiography of Julia''ను \citep{Reid1986}, అలాగే సమగ్ర స్మృతిరచనను \citep{Reid1996} ప్రచురించారు. రాబిన్సన్, హిల్బర్ట్ పదవ సమస్యపై చిన్న డాక్యుమెంటరీని జార్జ్ చిచ్చెరీ దర్శకత్వం వహించారు \citep{Csicsery2016}. రాబిన్సన్‌తో యూరీ మతియాసేవిచ్ సహకారం, ఆయన కృషిపై ఆమె ప్రభావం గురించి సంక్షిప్త స్మృతిరచనకు \citep{Matijasevich1992} చూడండి.
\end{reading}

\end{document}
